\section{Related work}
\label{sec:related}

\paragraph{Conformal prediction.}
Conformal prediction \citep{vovk2005algorithmic} gives finite-sample marginal coverage under exchangeability; the inductive (split) variant \citep{papadopoulos2002inductive,lei2018distribution} separates fitting from calibration, and cross-conformal methods such as the jackknife+ \citep{barber2021predictive} trade computation for sample efficiency.
Weighted conformal prediction handles known covariate shift \citep{tibshirani2019conformal}.

\paragraph{Set-valued classification.}
For classifier outputs $\hat p(x)\in\simplex$, the least-ambiguous-set score \citep{sadinle2019least}, adaptive prediction sets \citep{romano2020classification} and regularised variants \citep{angelopoulos2021uncertainty} differ in conditional behaviour, not in marginal validity.
Class-conditional and clustered calibration targets conditional coverage under many classes \citep{ding2023class}, and multivalid methods target coverage over rich groups \citep{jung2023batch}.
Credal-set predictors return sets of distributions \citep{javanmardi2024conformalized}.

\paragraph{Conditional coverage: limits and how it is measured.}
Exact distribution-free conditional coverage is impossible for non-atomic features, and approximate versions require relaxations \citep{vovk2012conditional,barber2021limits}; coverage conditional on the training set is a separate guarantee \citep{bian2023training}.
Practice therefore reports relaxed, \emph{estimated} summaries: the minimum coverage over feature slabs (worst-slab coverage, \citealp{cauchois2020knowing}), the maximal deviation of coverage across set-size strata \citep{angelopoulos2021uncertainty}, and the class-conditional coverage gap \citep{ding2023class}.
\citet{braun2025diagnostics} note in an appendix that the worst-slab estimate is systematically below the target even under exact conditional coverage, since the expectation of a minimum is at most the minimum of expectations, and propose classifier-based diagnostics.
In contrast, \citet{cauchois2020knowing} report in one multilabel simulation that the worst-slab ordering of methods agreed with the ordering by exact conditional coverage.
We take the first as established (Proposition~\ref{prop:min} makes it quantitative) and test the second: whether a method-dependent bias can change a ranking.

\paragraph{Compositional and multivariate responses.}
Compositional data analysis treats observations on the simplex through log-ratio geometry \citep{aitchison1982statistical,aitchison1986statistical}; zeros require special handling, and we use the geometries as score choices.
For vector-valued responses, conformal regions have been built from ellipsoidal scores based on residual covariances \citep{henderson2024ellipsoids}, from optimal-transport ranks of multivariate scores \citep{klein2025multivariate}, and, for outputs linked by linear constraints, from a projection step that provably shrinks split-conformal regions \citep{principato2024hierarchical}.
These works target $\mathbb R^d$ or linear structure; the non-negativity boundary and the choice of geometry are not their subject as far as the abstracts we read go.
Our search was limited and deliberately excluded work on conformal prediction with compositional responses (Appendix~\ref{app:repro}); we make no claim about its absence.

\paragraph{Gap.}
To our knowledge and within the sources consulted, three points are not addressed together: classification-set evaluation is defined for a discrete label, while here the response is continuous, so set size becomes a volume relative to an oracle region; the bias of the worst-cell estimator, and whether it can reverse a ranking, has not been quantified where the true conditional coverage is known; and the choice of geometry for a simplex-valued response is usually left to the analyst, although rank-based coverage is invariant to monotone score transformations but not to a change of geometry.
